% theme: inscribed_angle_mixed_problems
% 表：円周角の定理の混合計算15題。裏：等分された円周、弦の交点、円外の交点。
\MajorQuestionLesson{円周角の定理と角度}
\SetRowSpace{0mm}

\TypeQuestion{図の $x$ の大きさを求めなさい。点 $\mathrm{O}$ は円の中心、太線は直径である。同じ印の弧は等しく、弧のそばの数字は長さの比を表す。}{inscribed_angle_mixed}
\QQFigRowThree[32mm]{fig/lesson_02_mixed_01}{}{fig/lesson_02_mixed_02}{}{fig/lesson_02_mixed_03}{}
\QQFigRowThree[32mm]{fig/lesson_02_mixed_04}{}{fig/lesson_02_mixed_05}{}{fig/lesson_02_mixed_06}{}
\QQFigRowThree[32mm]{fig/lesson_02_mixed_07}{}{fig/lesson_02_mixed_08}{}{fig/lesson_02_mixed_09}{}
\QQFigRowThree[32mm]{fig/lesson_02_mixed_10}{}{fig/lesson_02_mixed_11}{}{fig/lesson_02_mixed_12}{}
\QQFigRowThree[32mm]{fig/lesson_02_mixed_13}{}{fig/lesson_02_mixed_14}{}{fig/lesson_02_mixed_15}{}

\AnswerColumnBreak
\LessonQuestionBreak
\SetRowSpace{1mm}

\TypeQuestion{円周を5等分する点を図のように $\mathrm{A}$ から $\mathrm{E}$ とし、$\mathrm{O}$ を円の中心とする。}{regular_pentagon_angles}
\FigProblemBlock{40mm}{fig/lesson_02_equal_05}{
\QQ{$\Ang{AOC}$ の大きさを求めなさい。}{}
\QQ{$\Ang{ACB}$ の大きさを求めなさい。}{}
\QQ{$\Ang{BAE}$ の大きさを求めなさい。}{}
\QQ{$\Ang{AFB}$ の大きさを求めなさい。}{}
}{}

\TypeQuestion{円周を6等分する点を図のように $\mathrm{A}$ から $\mathrm{F}$ とし、$\mathrm{O}$ を円の中心とする。}{regular_hexagon_angles}
\FigProblemBlock{40mm}{fig/lesson_02_equal_06}{
\QQ{$\Ang{AOC}$ の大きさを求めなさい。}{}
\QQ{$\Ang{BAC}$ の大きさを求めなさい。}{}
\QQ{$\Ang{ABD}$ の大きさを求めなさい。}{}
\QQ{$\Ang{AGB}$ の大きさを求めなさい。}{}
}{}

\TypeQuestion{円周を8等分する点を図のように $\mathrm{A}$ から $\mathrm{H}$ とする。}{regular_octagon_angles}
\FigProblemBlock{40mm}{fig/lesson_02_equal_08}{
\QQ{$\Ang{AIB}$ の大きさを求めなさい。}{}
\QQ{$\Ang{ABE}$ の大きさを求めなさい。}{}
\QQ{$\Ang{EAB}$ の大きさを求めなさい。}{}
\QQ{$\Ang{DAE}$ の大きさを求めなさい。}{}
}{}

\TypeQuestion{円周を9等分する点を図のように $\mathrm{A}$ から $\mathrm{I}$ とする。}{regular_hexagon_exterior_angles}
\FigProblemBlock{40mm}{fig/lesson_02_equal_external}{
\QQ{$\Ang{CBD}$ の大きさを求めなさい。}{}
\QQ{$\Ang{BGH}$ の大きさを求めなさい。}{}
\QQ{$\Ang{BJH}$ の大きさを求めなさい。}{}
\QQ{$\Ang{BKH}$ の大きさを求めなさい。}{}
}{}
